\chapter{Conclusões e Trabalhos Futuros}
\label{chap:conclusoes-e-trabalhos-futuros}

Este trabalho apresentou o EDGAR, uma arquitetura de Geração Aumentada por Recuperação Multimodal (MRAG) em granularidade fina voltada à consulta de materiais didáticos em PDF. A implementação foi organizada em dois fluxos complementares, combinando a detecção de componentes documentais, a extração textual, a descrição de elementos visuais por modelos de visão e linguagem e a indexação vetorial no Qdrant. A associação de coordenadas geométricas aos registros permitiu fundamentar as respostas geradas pelo modelo de linguagem e destacar a região de origem das evidências na interface \textit{web}.

Dentre as dificuldades encontradas, destaca-se a limitação de créditos computacionais na plataforma em nuvem Modal, que restringiu o processamento a 500 das 578 páginas planejadas e reduziu a avaliação final para 21 consultas. Além disso, a ausência de transcrição específica para fórmulas matemáticas isoladas e a ocorrência de falhas pontuais de leitura pelo modelo de visão impactaram o desempenho em consultas dependentes de equações.

Os resultados obtidos evidenciam a viabilidade técnica da recuperação em nível de componentes associada ao enriquecimento visual, que alcançou a maior média na avaliação por critérios de qualidade das respostas entre as configurações comparadas. A abordagem completa do EDGAR mostrou-se capaz de elevar o desempenho nas consultas dependentes de figuras e de raciocínio multimodal, além de assegurar a rastreabilidade visual das evidências no próprio documento.

Para trabalhos futuros, diversas melhorias podem ser exploradas. Uma possível linha de pesquisa envolve a transcrição específica de fórmulas matemáticas isoladas para representação em LaTeX, bem como o envio direto dos recortes visuais ao modelo gerador durante a produção da resposta. Outra direção consiste em avaliar estratégias de indexação híbrida e de reordenação semântica das evidências recuperadas. A ampliação do conjunto de documentos e a realização de avaliações com especialistas humanos e estudantes também podem contribuir para validar a utilidade prática da ferramenta em cenários reais de ensino.
